% ============================================================
% Chapter 8: Sustainability and Ethical Implications
%   Section: Ethics of Data Management and Analysis
%     Subsection: Collection of Third-Party Content
% Included from chapters/08-sustainability-ethics/02-data-ethics/data-ethics.tex
% ============================================================

\subsection{Collection of Third-Party Content}
\label{subsec:third-party-content}

The system reads articles that other people wrote and published. What
it does with them is limited in four ways.

\begin{itemize}
    \item \textbf{What is stored.} The text of an analysed article is
    kept in the data lake, which is local and not published, so that a
    verdict can be traced back to what was read. Evidence pages are
    read for the claim they bear on; only their title, a quoted passage
    and their address are kept with the verdict, and the evaluation
    harness stores no page body at all. Raw evaluation runs, which hold
    third-party titles and quotes, are not committed to the repository.
    \item \textbf{What is shown.} The public reader view shows an
    article's lead and a link to the original, never its body or a
    scraped page; a test holds that rule. The original publisher keeps
    the reader.
    \item \textbf{How it asks.} The scraper names itself in its
    user-agent, asks each source for a few articles per run rather than
    crawling it, and is run on request, never on a timer. It does not
    read \texttt{robots.txt} (Section~\ref{subsec:scraping-limitations}),
    which is a limitation to correct before any public deployment.
    \item \textbf{When it is refused.} A source that refuses automated
    access is disabled, not worked around: EFE and Reuters, which
    answered every request with 403 and 401 for two weeks, were disabled
    on 6~October. The second search route, when challenged as a bot,
    pauses rather than pretending to be a browser.
\end{itemize}

Fact-checking organisations' pages are used as evidence when a search
returns them, and the two Spanish organisations are rated as sources,
but they are not ingested: a fact-check quotes the claim it debunks, and
the system would otherwise present that quotation as an article's claim.
